\documentclass[tikz,border=8pt]{standalone}
\usepackage{tikz}
\usepackage{amsmath}
\usepackage{amssymb}
\usepackage{pgfplots}
\pgfplotsset{compat=1.18}
\usepackage{ctex}
\usetikzlibrary{calc,positioning,arrows.meta,matrix}

% LC 70 爬楼梯 & LC 198 打家劫舍
% 爬楼梯：dp[i] = dp[i-1] + dp[i-2]
% 打家劫舍：dp[i] = max(dp[i-1], dp[i-2] + nums[i])

\begin{document}
\begin{tikzpicture}[
  every node/.style={font=\small},
  cell/.style={draw=gray!70, minimum size=0.8cm, inner sep=0pt, thick, fill=white},
  hlcell/.style={draw=blue!60, minimum size=0.8cm, inner sep=0pt, thick, fill=blue!8},
  curcell/.style={draw=orange!70, minimum size=0.8cm, inner sep=0pt, very thick, fill=orange!8},
  arr/.style={->,>=Stealth,thick,blue!70},
  annbox/.style={draw=gray!50, fill=white, rounded corners=3pt, inner sep=5pt},
]

%% ── 标题 ────────────────────────────────────────────────────
\node[font=\large\bfseries] at (5.5, 0.9)
  {一维 DP 状态转移（LC 70 爬楼梯 / LC 198 打家劫舍）};
\node[font=\small, gray] at (5.5, 0.3)
  {横向数组展示 dp[i] 初始化与逐步转移过程};

%% ════════════════════════════════════════════════════════════
%% 上半：LC 70 爬楼梯
%% dp[0]=1, dp[1]=1, dp[i] = dp[i-1] + dp[i-2]  (n=6)
%% ════════════════════════════════════════════════════════════

\node[font=\normalsize\bfseries] at (5.5, -0.5)
  {LC 70：爬楼梯（n=6，求到达第6阶的方法数）};

% 初始化状态标签
\node[font=\scriptsize, gray] at (-0.7, -1.3) {索引};
\node[font=\scriptsize, gray] at (-0.7, -2.1) {dp[i]};

% 索引行
\foreach \i in {0,...,6} {
  \node[font=\scriptsize, gray] at (\i*1.1+0.55, -1.3) {\i};
}

% dp 数组格子（完整填充后）
\node[hlcell] at (0.55, -2.1) {1};
\node[hlcell] at (1.65, -2.1) {1};
\node[cell]   at (2.75, -2.1) {2};
\node[cell]   at (3.85, -2.1) {3};
\node[cell]   at (4.95, -2.1) {5};
\node[cell]   at (6.05, -2.1) {8};
\node[curcell]at (7.15, -2.1) {13};

% dp 值标签
\node[font=\tiny, blue!70, above=2pt] at (0.55, -1.75) {初始};
\node[font=\tiny, blue!70, above=2pt] at (1.65, -1.75) {初始};
\node[font=\tiny, orange!70, above=2pt] at (7.15, -1.75) {目标};

% 转移箭头：dp[5]+dp[6] → dp[7]（示意）
\draw[arr] (1.65, -2.55) to[out=-60,in=-120] (2.75, -2.55);
\draw[arr] (0.55, -2.55) to[out=-70,in=-150] (2.75, -2.55);
\node[font=\tiny, gray] at (1.65, -3.15) {dp[1]$+$dp[0]};

\draw[arr] (2.75, -2.55) to[out=-60,in=-120] (3.85, -2.55);
\draw[arr] (1.65, -2.55) to[out=-70,in=-150] (3.85, -2.55);
\node[font=\tiny, gray] at (3.3, -3.15) {dp[2]$+$dp[1]};

\draw[arr] (5.05, -2.55) to[out=-60,in=-120] (6.05, -2.55);
\draw[arr] (3.95, -2.55) to[out=-60,in=-150] (6.05, -2.55);

\draw[arr,orange!70,very thick] (6.05, -2.55) to[out=-60,in=-120] (7.15, -2.55);
\draw[arr,orange!70,very thick] (4.95, -2.55) to[out=-70,in=-150] (7.15, -2.55);

% 转移公式框
\node[annbox, font=\small, align=left] at (10.0, -2.1)
  {\textbf{转移公式}\\
   $\text{dp}[0]=1,\; \text{dp}[1]=1$\\
   $\text{dp}[i] = \text{dp}[i-1] + \text{dp}[i-2]$\\
   $\Rightarrow$ 方法数 $= \textbf{13}$};

%% ════════════════════════════════════════════════════════════
%% 中间：LC 198 打家劫舍
%% nums = [2, 7, 9, 3, 1]
%% dp[0]=2, dp[1]=max(2,7)=7, dp[i]=max(dp[i-1], dp[i-2]+nums[i])
%% ════════════════════════════════════════════════════════════

\node[font=\normalsize\bfseries] at (5.5, -5.0)
  {LC 198：打家劫舍（nums = [2, 7, 9, 3, 1]）};

% nums 标签行
\node[font=\scriptsize, gray] at (-0.7, -5.8) {索引};
\node[font=\scriptsize, gray] at (-0.7, -6.6) {nums};
\node[font=\scriptsize, gray] at (-0.7, -7.4) {dp[i]};

\foreach \i/\v in {0/2, 1/7, 2/9, 3/3, 4/1} {
  \node[font=\scriptsize, gray] at (\i*1.3+0.65, -5.8) {\i};
  % nums 格子（白底灰边）
  \node[cell] at (\i*1.3+0.65, -6.6) {\v};
}

% dp 数组：dp=[2,7,11,11,12]
\node[hlcell] at (0.65, -7.4) {2};
\node[hlcell] at (1.95, -7.4) {7};
\node[cell]   at (3.25, -7.4) {11};
\node[cell]   at (4.55, -7.4) {11};
\node[curcell]at (5.85, -7.4) {12};

\node[font=\tiny, blue!70, above=2pt] at (0.65, -7.05) {初始};
\node[font=\tiny, blue!70, above=2pt] at (1.95, -7.05) {初始};
\node[font=\tiny, orange!70, above=2pt] at (5.85, -7.05) {最优};

% 分步演化（Step 2: dp[2]）
\draw[arr, blue!70] (0.65, -7.85) to[out=-60,in=-120] (3.25, -7.85);
\node[font=\tiny, gray, below] at (1.95, -8.1) {dp[0]+nums[2]=11};
\draw[arr, gray!60] (1.95, -7.85) to[out=-60,in=-120] (3.25, -7.85);

% Step 3: dp[3]
\draw[arr, gray!60] (3.25, -7.85) to[out=-60,in=-120] (4.55, -7.85);
\draw[arr, gray!60] (1.95, -7.85) to[out=-70,in=-150] (4.55, -7.85);

% Step 4: dp[4] = max(dp[3], dp[2]+nums[4]) = max(11, 12) = 12
\draw[arr, orange!70, very thick] (3.25, -7.85) to[out=-60,in=-120] (5.85, -7.85);
\draw[arr, gray!50] (4.55, -7.85) to[out=-60,in=-150] (5.85, -7.85);
\node[font=\tiny, orange!70, below] at (4.55, -8.35)
  {max(11, 11+1=12)=\textbf{12}};

% 转移公式框
\node[annbox, font=\small, align=left] at (10.0, -7.4)
  {\textbf{转移公式}\\
   $\text{dp}[0]=\text{nums}[0]$\\
   $\text{dp}[1]=\max(\text{nums}[0],\text{nums}[1])$\\
   $\text{dp}[i]=\max(\text{dp}[i\!-\!1],\; \text{dp}[i\!-\!2]+\text{nums}[i])$\\
   $\Rightarrow$ 最大金额 $= \textbf{12}$};

%% ── 分隔线 ──────────────────────────────────────────────────
\draw[gray!40, dashed] (-1.2, -4.4) -- (12.2, -4.4);

%% ── 图例 ──────────────────────────────────────────────────
\node[annbox, font=\scriptsize, align=left] at (5.5, -10.2)
  {\textbf{图例：}\ 蓝色格子 = 初始化值 \quad 橙色格子 = 最终目标值
   \quad 蓝色箭头 = 关键转移路径 \quad 灰色箭头 = 其余转移};

\end{tikzpicture}
\end{document}
